% from ega4-errata-addenda-list3-fr.tex
\oldpage[IV]{349}
\begin{remark}[3.1.10.5]
\label{IV.3.1.10.5}
On notera que même lorsque $f$ est un morphisme propre, on n'a pas
nécessairement
\[
 f(\operatorname{Ass}(\sh{F}))\subset
 \operatorname{Ass}(f_*(\sh{F})).
\]
On le voit en prenant pour $Y$ le spectre d'un corps et pour $X$ un
espace projectif sur $k$, par exemple $\mathbf P^1_k$; alors
$\sh{F}=\sh{O}_X(-1)$ est un $\sh{O}_X$-Module cohérent $\neq0$ tel
que $\Gamma(X,\sh{F})=0$ (III, 2.1.13), donc
$\operatorname{Ass}(f_*(\sh{F}))=\varnothing$ et
$\operatorname{Ass}(\sh{F})\neq\varnothing$.
\end{remark}

\paragraph{(Err\textsubscript{IV}, 18).}
Après (IV, 3.2.8), insérer après la ligne 17 de la p.~43 :

\begin{remark}[3.2.9]
\label{IV.3.2.9}
Pour tout point $s$ d'un préschéma $X$, notons
$i_s:\operatorname{Spec}(k(s))\to X$ le morphisme canonique, qui est
un monomorphisme quasi-compact. Pour tout $k(s)$-module $G(s)$,
considéré comme Module sur $\operatorname{Spec}(k(s))$,
$(i_s)_*(G(s))$ est donc un $\sh{O}_X$-Module quasi-cohérent
{\hyperref[canon-index.ega4.c04.statement.1-7-4]{(1.7.4)}}. On vérifie aussitôt
que la fibre de $(i_s)_*(G(s))$ est nulle en tout point
$x\notin\overline{\{s\}}$ et isomorphe à $G(s)$ dans le cas contraire.
Si $X$ est localement noethérien et $G(s)$ de longueur finie, il est
clair que $\operatorname{Ass}((i_s)_*(G(s)))$ est réduit au seul point
$s$, autrement dit $(i_s)_*(G(s))$ est irredondant.

Cela étant, supposons $X$ localement noethérien, et soit $\sh{F}$ un
$\sh{O}_X$-Module cohérent. Soit d'autre part $S$ une partie discrète
de $X$, et pour tout $s\in S$, posons
$\sh{F}(s)=i_s^*(\sh{F})$, qui est un $k(s)$-module; considérons pour
chaque $s\in S$ un module quotient $G(s)$ de $\sh{F}(s)$ de longueur
finie. On a donc un homomorphisme canonique
\[
 u:\sh{F}\longrightarrow
 \bigoplus_{s\in S}(i_s)_*(\sh{F}(s))\longrightarrow
 \bigoplus_{s\in S}(i_s)_*(G(s))
\]
de $\sh{O}_X$-Modules cohérents (puisque $S$, étant discret, est
localement fini dans $X$). Si $\sh{H}=\operatorname{Ker}u$,
$\sh{G}=\sh{F}/\sh{H}$ est un $\sh{O}_X$-Module cohérent et on déduit
de $u$ un homomorphisme injectif
\[
 v:\sh{G}\longrightarrow\bigoplus_{s\in S}(i_s)_*(G(s)).
\]
Si $\sh{G}^{(s)}$ est l'image de $\sh{G}$ dans $(i_s)_*(G(s))$,
$(\sh{G}^{(s)})_{s\in S}$ est une décomposition irredondante, qui est
réduite lorsque les $G(s)$ sont tous $\neq0$, et telle que
$\operatorname{Ass}(\sh{G})=S$ ($\sh{G}$ étant donc sans cycle premier
associé immergé). On a ainsi obtenu une bijection
$(G(s))_{s\in S}\leftrightarrow\sh{G}$, où $S$ parcourt l'ensemble des
parties discrètes de $X$ pour lesquelles $\sh{F}(s)\neq0$ pour tout
$s\in S$, et, pour chaque $S$, $(G(s))_{s\in S}$ parcourt l'ensemble
des quotients $\neq0$ et de longueur finie des $\sh{F}(s)$; d'autre
part $\sh{G}$ parcourt l'ensemble des $\sh{O}_X$-Modules quotients
cohérents de $\sh{F}$, sans cycle premier immergé. La bijection
réciproque associe à chaque $\sh{G}$ sa décomposition irredondante
réduite.
\end{remark}

\paragraph{(Err\textsubscript{IV}, 19).}
Dans (IV, 4.2.1), avant la ligne 17 du bas de la p.~55, insérer :

\begin{remark}[4.2.1.4]
\label{IV.4.2.1.4}
Sous les hypothèses de {\hyperref[canon-index.ega4.c04.prop.4-2-1]{(4.2.1)}}, on a
en outre
\begin{equation}
 \dim(K\otimes_kL)=
 \inf(\operatorname{deg.tr.}_kK,\operatorname{deg.tr.}_kL).
 \label{IV.4.2.1.5}\tag{4.2.1.5}
\end{equation}
où il doit être entendu que si les deux corps $K$, $L$ sont de degré de
transcendance infini (qui peuvent être des cardinaux quelconques), le
premier membre doit être pris égal à $+\infty$ (la dimension considérée
ici est celle définie dans ($0$, 16.1.1)).

Supposons en effet par exemple que
$\operatorname{deg.tr.}_kK\leq\operatorname{deg.tr.}_kL$ et en premier
lieu que $\operatorname{deg.tr.}_kK=n<+\infty$. On a vu dans la preuve
de {\hyperref[canon-index.ega4.c04.prop.4-2-1]{(4.2.1)}} que $K\otimes_kL$ est une
extension algébrique de $K'\otimes_kL$, et a donc même dimension
($0$, 16.1.5). On peut donc se borner au cas où
$K=k[T_1,\ldots,T_n]$; alors $K\otimes_kL$ est un anneau de fractions
de
\[
 k[T_1,\ldots,T_n]\otimes_kL=L[T_1,\ldots,T_n],
\]
qui est de dimension $n$ {\hyperref[canon-index.ega4.c04.cor.5-5-4]{(5.5.4)}}, et
par suite $\dim(K\otimes_kL)\leq n$ ($0$, 16.1.3.1). Pour prouver que
$\dim(K\otimes_kL)=n$, il suffit évidemment de montrer qu'il y a dans
$\operatorname{Spec}(L[T_1,\ldots,T_n])$ un point rationnel sur $L$
qui appartient à $\operatorname{Spec}(K\otimes_kL)$
{\hyperref[canon-index.ega4.c04.cor.5-2-3]{(5.2.3)}}, ou encore (I, 1.7.4) qu'il
existe un $L$-homomorphisme $u:K\otimes_kL\to L$. Mais par définition
du produit tensoriel, cela revient à prouver l'existence d'un
$k$-homomorphisme $v:K\to L$, ce qui est évident en vertu de
l'hypothèse
$\operatorname{deg.tr.}_kK\leq\operatorname{deg.tr.}_kL$.

Si $\operatorname{deg.tr.}_kK\leq\operatorname{deg.tr.}_kL$ et si ces
deux cardinaux sont infinis, on voit de la même façon qu'il y a un point
rationnel dans $\operatorname{Spec}(L[(t_\alpha)_{\alpha\in I}])$ qui
appartient à $\operatorname{Spec}(K\otimes_kL)$; mais le noyau de tout
$L$-homomorphisme de $L[(t_\alpha)_{\alpha\in I}]$ dans $L$ est un
idéal maximal engendré par une famille de la forme
$(t_\alpha-a_\alpha)_{\alpha\in I}$ avec $a_\alpha\in L$. Il est clair
qu'un tel idéal maximal contient une suite infinie décroissante
d'idéaux premiers distincts, donc l'anneau local de
$L[(t_\alpha)_{\alpha\in I}]$ correspondant est non noethérien et de
dimension infinie; comme c'est un anneau local de $K\otimes_kL$, on
voit que ce dernier est non noethérien et de dimension infinie.
\end{remark}

